% control_architecture.tex — Electronic control system block diagram.
% \input inside a figure environment. Requires KJO_tikz_styles.tex.

\begin{tikzpicture}[x=1cm, y=1cm, node distance=1.2cm]

%── RCU (left) ────────────────────────────────────────────────────────────────
\node[rectangle, draw=KJOblue, line width=1.5pt, fill=KJOblue!8,
      rounded corners=6pt, minimum width=4.0cm, minimum height=5.5cm,
      font=\small\bfseries, align=center, label={[font=\small\bfseries,
      text=KJOblue]above:Remote Control Unit (RCU)}]
  (rcu) at (0, 0) {};

\node[mcublock, minimum width=3.5cm] (rcu_mcu) at (0, 1.8) {Feather M0};
\node[hwblock, minimum width=3.5cm] (rcu_disp) at (0, 0.7) {3.5\textquotedbl\ TFT\\+ Touchscreen};
\node[hwblock, minimum width=3.5cm] (rcu_rad) at (0, -0.4) {RFM95 LoRa};
\node[hwblock, minimum width=3.5cm] (rcu_sd) at (0, -1.5) {SD + RTC};
\node[powerblock, minimum width=3.5cm] (rcu_bat) at (0, -2.4) {LiPo Battery};

%── EMU (center) ──────────────────────────────────────────────────────────────
\node[rectangle, draw=KJOred, line width=1.5pt, fill=KJOred!5,
      rounded corners=6pt, minimum width=4.2cm, minimum height=8.0cm,
      font=\small\bfseries, align=center, label={[font=\small\bfseries,
      text=KJOred]above:Engine Management Unit (EMU)}]
  (emu) at (6.5, -1.3) {};

\node[mcublock, minimum width=3.7cm] (emu_mcu) at (6.5, 2.0) {Feather M0};
\node[hwblock, minimum width=3.7cm] (emu_rad) at (6.5, 1.0) {RFM95 LoRa};
\node[hwblock, minimum width=3.7cm] (emu_oled) at (6.5, 0.0) {SH1107 OLED};
\node[hwblock, minimum width=3.7cm] (emu_gpio) at (6.5, -1.0) {MCP23X17 GPIO\\Expander};
\node[hwblock, minimum width=3.7cm] (emu_pwm) at (6.5, -2.0) {PWM Servo Wing\\(Ox, Fuel)};
\node[hwblock, minimum width=3.7cm] (emu_adc) at (6.5, -3.0) {ADS1015 ADC\\+ MAX31856 TC};
\node[hwblock, minimum width=3.7cm] (emu_sd) at (6.5, -4.0) {SD + RTC};
\node[powerblock, minimum width=3.7cm] (emu_bat) at (6.5, -4.8) {LiPo Battery};

%── GSEMU (right of EMU) ──────────────────────────────────────────────────────
\node[rectangle, draw=KJOgreen, line width=1.5pt, fill=KJOgreen!5,
      rounded corners=6pt, minimum width=4.0cm, minimum height=5.5cm,
      font=\small\bfseries, align=center, label={[font=\small\bfseries,
      text=KJOgreen]above:GSE Mgmt Unit (GSEMU)}]
  (gsemu) at (13.0, 0) {};

\node[mcublock, minimum width=3.5cm] (gse_mcu) at (13.0, 1.8) {Feather M0};
\node[hwblock, minimum width=3.5cm] (gse_oled) at (13.0, 0.7) {SH1107 OLED};
\node[hwblock, minimum width=3.5cm] (gse_pwm) at (13.0, -0.4) {PWM Servo Wing\\(Fill Valve)};
\node[hwblock, minimum width=3.5cm] (gse_sd) at (13.0, -1.5) {SD + RTC};
\node[powerblock, minimum width=3.5cm] (gse_bat) at (13.0, -2.4) {LiPo Battery};

%── LCMU (right of GSEMU) ─────────────────────────────────────────────────────
\node[rectangle, draw=KJOpurple, line width=1.5pt, fill=KJOpurple!5,
      rounded corners=6pt, minimum width=4.0cm, minimum height=5.5cm,
      font=\small\bfseries, align=center, label={[font=\small\bfseries,
      text=KJOpurple]above:Load Cell Mgmt Unit (LCMU)}]
  (lcmu) at (19.5, 0) {};

\node[mcublock, minimum width=3.5cm] (lcmu_mcu)  at (19.5, 1.8) {Feather RP2040\\Adalogger};
\node[hwblock, minimum width=3.5cm]  (lcmu_oled) at (19.5, 0.6) {SH1107 OLED};
\node[hwblock, minimum width=3.5cm]  (lcmu_cell) at (19.5, -0.5) {Mux + 3$\times$NAU7802\\(Load Cells)};
\node[hwblock, minimum width=3.5cm]  (lcmu_sd)   at (19.5, -1.6) {SD + RTC};
\node[powerblock, minimum width=3.5cm] (lcmu_bat) at (19.5, -2.4) {LiPo Battery};

%── LoRa Link (RCU ↔ EMU) ─────────────────────────────────────────────────────
\draw[radioline]
  (rcu.east |- rcu_rad) -- (emu.west |- emu_rad)
  node[midway, above, font=\footnotesize\itshape, text=KJOblue!70]
    {900 MHz LoRa};
\node[font=\tiny, text=KJOgray, align=center]
  at (3.25, 0.55) {Commands\\Responses};

%── CAN Bus (EMU -- GSEMU -- LCMU) ────────────────────────────────────────────
% Shared bus, not point-to-point: replaces the retired CMD/STATUS GPIO
% handshake for Fill-valve coordination, and is LCMU's only link to the
% other two units. Drawn as one bus line with tap stubs, matching CAN's
% actual multi-drop topology rather than the old wired handshake's
% point-to-point arrows.
% Routed below all three outer boxes (clear of every inner node) rather
% than through their interior -- GSEMU/LCMU tap in via a short vertical
% stub from their box's south edge; EMU's box already extends past this
% height, so its east-edge border point is used directly, no stub needed.
\coordinate (canbus_y) at (0, -3.3);
\draw[busline, -, line width=2pt]
  (emu.east |- canbus_y) -- (lcmu.east |- canbus_y)
  node[midway, below, font=\tiny, text=KJOgray] {CAN bus (1 Mbps)};
\draw[busline, -] (gsemu.south) -- (gsemu.south |- canbus_y);
\draw[busline, -] (lcmu.south)  -- (lcmu.south  |- canbus_y);

%── Umbilical disconnect-sense (EMU <-> GSEMU only, separate from CAN) ────────
\draw[gpioline, draw=KJOgray!70, -, dashed]
  ([yshift=-14pt] emu.east |- emu_gpio) --
  ([yshift=-14pt] gsemu.west |- emu_gpio)
  node[midway, below, font=\tiny, text=KJOgray!70!black] {disconnect sense (ground-ref.)};

%── Physical separation note ──────────────────────────────────────────────────
\node[font=\footnotesize\itshape, text=KJOgray, align=center]
  at (3.25, -5.2) {Operator\\(remote)};
\node[font=\footnotesize\itshape, text=KJOgray, align=center]
  at (6.5, -5.85) {In-rocket};
\node[font=\footnotesize\itshape, text=KJOgray, align=center]
  at (13.0, -4.3) {Launch pad\\(ground)};
\node[font=\footnotesize\itshape, text=KJOgray, align=center]
  at (19.5, -4.3) {Test stand};

%── Ground line ───────────────────────────────────────────────────────────────
\draw[KJOgray!40, line width=0.8pt, dashed]
  (0, -5.0) -- (6.5, -5.0)
  node[midway, below, font=\tiny, text=KJOgray] {radio link};

\end{tikzpicture}
